\begin{description}
\item[(a)] If $a$ is the semi-major axis, $M_S$ the mass of the Sun and $G$ the gravitational constant,
\[
v_{\max}=\sqrt{\frac{GM_S}{a}\frac{1+e}{1-e}},\qquad
v_{\min}=\sqrt{\frac{GM_S}{a}\frac{1-e}{1+e}}.
\]\hfill(1.5 points)
Then $k=v_{\max}/v_{\min}=(1+e)/(1-e)$ and $k(1-e)=1+e$, so
\[
e=\frac{k-1}{k+1}=\frac{v_{\max}-v_{\min}}{v_{\max}+v_{\min}}.
\]\hfill(0.5 point)
Numerically, $e=(19-14)/(19+14)\approx0.15$. \hfill(1 point)
\item[(b)] At distance $r$ from the Sun,
\[
v(r)=\sqrt{GM_S\left(\frac2r-\frac1a\right)}.
\]\hfill(1 point)
At either co-vertex, the asteroid is at the intersection between the ellipse and its semi-minor axis, so $r=a$. \hfill(1 point)
Thus
\[
v(a)=\sqrt{\frac{GM_S}{a}}=\sqrt{v_{\max}v_{\min}}\approx16.31\,\mathrm{km/s}.
\]\hfill(1 point)
\item[(c)] For a comet from the outer edge of the Solar System, the orbit can be assumed parabolic. Its total energy is zero, so
\[
v=\sqrt{\frac{2GM_\odot}{d}},\qquad d_{\min}=\frac{2GM_\odot}{v_{\max}^2}.
\]\hfill(1.5 points)
Numerically,
\[
d_{\min}=\frac{2\times6.673\times10^{-11}\times1.989\times10^{30}}{(9\times10^4)^2}\approx0.22\,\mathrm{au}.
\]\hfill(0.5 point)
For the asteroid, $a=3.34\,\mathrm{au}$ and $r_{\max}=a(1+e)=3.84\,\mathrm{au}$. \hfill(1.5 points)
The comet--asteroid distance is $D=d_{\min}+r_{\max}=0.22+3.84=4.06\,\mathrm{au}$. \hfill(0.5 point)
\end{description}
